\documentclass[UTF8]{ctexart}
\usepackage{amsmath, amssymb, geometry}
\usepackage{xcolor}
\usepackage{enumitem}
\usepackage{tcolorbox}
\usepackage{cases}
\usepackage{fancyhdr}
\geometry{a4paper, margin=1in}

% 定义红色环境
\newenvironment{redenv}{\color{red}}{}

% 定义详细解析环境
\newenvironment{detailedanalysis}{%
    \color{red}
    \begin{enumerate}[label=\textbf{第\arabic*步：}, leftmargin=*, align=left, itemsep=1.5ex]
}{%
    \end{enumerate}
}

% 定义概念框
\newenvironment{conceptbox}{%
    \begin{tcolorbox}[colback=red!5, colframe=red!50!black, title=概念解析, fonttitle=\bfseries]
}{%
    \end{tcolorbox}
}

% 定义公式推导环境
\newenvironment{formuladerivation}{%
    \begin{enumerate}[label={\textbf{公式\arabic*：}}, leftmargin=*, align=left]
}{%
    \end{enumerate}
}

% 定义题目和解析的分隔线
\newcommand{\problemrule}{\noindent\rule{\textwidth}{1.5pt}\vspace{-0.8em}}
\newcommand{\analysisrule}{\noindent\rule{\textwidth}{0.8pt}\vspace{-0.8em}}

% 宏定义：方便排版选择题选项
\newcommand{\choicesFour}[4]{
    \begin{enumerate}[label=\Alph*., itemsep=0pt, parsep=0pt, topsep=5pt, labelsep=0.5em]
        \begin{minipage}{0.22\linewidth} \item #1 \end{minipage}
        \begin{minipage}{0.22\linewidth} \item #2 \end{minipage}
        \begin{minipage}{0.22\linewidth} \item #3 \end{minipage}
        \begin{minipage}{0.22\linewidth} \item #4 \end{minipage}
    \end{enumerate}
}
\newcommand{\choicesTwo}[4]{
    \begin{enumerate}[label=\Alph*., itemsep=0pt, parsep=0pt, topsep=5pt, labelsep=0.5em]
        \begin{minipage}{0.45\linewidth} \item #1 \end{minipage}
        \begin{minipage}{0.45\linewidth} \item #2 \end{minipage}
        \begin{minipage}{0.45\linewidth} \item #3 \end{minipage}
        \begin{minipage}{0.45\linewidth} \item #4 \end{minipage}
    \end{enumerate}
}
\newcommand{\choices}[4]{
    \begin{enumerate}[label=\Alph*., itemsep=0pt, parsep=0pt, topsep=5pt, labelsep=0.5em]
        \item #1
        \item #2
        \item #3
        \item #4
    \end{enumerate}
}

\newcommand{\blank}[1]{\underline{\hspace{#1}}}

\begin{document}

\section*{第三章 连续（提高）}

% ========== 第1题 ==========
\problemrule
\section*{【题目1】间断点类型}
已知函数 $ f(x) = \frac{x+2}{x^{2}-x} $ ，则 $ x = 0 $ 是函数 $ f(x) $ 的 \blank{1cm}。
\choicesFour{可去间断点}{跳跃间断点}{无穷间断点}{连续点}

\begin{redenv}
\analysisrule
\subsection*{逐步解析}

\begin{detailedanalysis}
\item \textbf{问题本质分析}
判别函数间断点的类型，关键在于计算该点的左右极限。

\item \textbf{详细求解过程}
\begin{enumerate}[label=(\arabic*)]
\item \textbf{确定定义域}：分母 $x^2 - x \neq 0 \implies x(x-1) \neq 0 \implies x \neq 0$ 且 $x \neq 1$。
所以 $x=0$ 是间断点。
\item \textbf{计算极限}：
$$ \lim_{x \to 0} f(x) = \lim_{x \to 0} \frac{x+2}{x(x-1)} $$
当 $x \to 0$ 时，分子 $x+2 \to 2$，分母 $x(x-1) \to 0 \cdot (-1) = 0$。
所以极限为 $\infty$。
\item \textbf{结论}：若函数在某点的极限为无穷大，则该点为无穷间断点。
\end{enumerate}

\item \textbf{最终答案}
\textbf{答案}：C
\end{detailedanalysis}

\subsection*{考点深度分析}
\begin{conceptbox}
\textbf{间断点分类}：
1. 第一类：左右极限均存在。
   - 可去：左右极限相等。
   - 跳跃：左右极限不相等。
2. 第二类：左右极限至少有一个不存在。
   - 无穷：极限为无穷大。
   - 振荡：极限振荡不存在。
\end{conceptbox}
\end{redenv}

% ========== 第2题 ==========
\problemrule
\section*{【题目2】无穷大与无穷小}
当 $ x \rightarrow +\infty $ 时， $ x\sin x $ 是 \blank{1cm}。
\choicesFour{无穷大量}{无穷小量}{无界变量}{有界变量}

\begin{redenv}
\analysisrule
\subsection*{逐步解析}

\begin{detailedanalysis}
\item \textbf{详细求解过程}
\begin{enumerate}[label=(\arabic*)]
\item \textbf{分析无穷大量}：若 $\lim_{x \to +\infty} |f(x)| = +\infty$，则为无穷大量。
取 $x_n = 2n\pi + \frac{\pi}{2}$，则 $x_n \sin x_n = x_n \to +\infty$。
取 $x_n = n\pi$，则 $x_n \sin x_n = 0$。
函数值在有些点很大，有些点为0，所以极限不是无穷大。
\item \textbf{分析有界性}：
由于 $x \to +\infty$ 可以取任意大的值，且 $\sin x$ 周期性取 $\pm 1$，所以 $x \sin x$ 可以取任意大的值（无界）。
例如取 $x_k = 2k\pi + \frac{\pi}{2}$，则 $f(x_k) \to +\infty$。
\item \textbf{结论}：既不是无穷大量（因为有零点序列），也不是有界变量（因为有无穷大序列）。称为无界变量。
\end{enumerate}

\item \textbf{最终答案}
\textbf{答案}：C
\end{detailedanalysis}
\end{redenv}

% ========== 第3题 ==========
\problemrule
\section*{【题目3】无穷小阶的比较}
当 $ x \rightarrow 0 $ 时， $ x^{2}\sin\frac{1}{x} $ 是 \blank{1cm}。
\choicesTwo
{与 $x$ 等价无穷小}
{$x$ 的低阶无穷小}
{与 $x$ 等价无穷小}
{$x$ 的高阶无穷小}

\begin{redenv}
\analysisrule
\subsection*{逐步解析}

\begin{detailedanalysis}
\item \textbf{详细求解过程}
比较 $f(x) = x^2 \sin \frac{1}{x}$ 与 $g(x) = x$ 的阶。
计算极限：
$$ \lim_{x \to 0} \frac{x^2 \sin \frac{1}{x}}{x} = \lim_{x \to 0} x \sin \frac{1}{x} $$
因为 $x$ 是无穷小量，$\sin \frac{1}{x}$ 是有界变量。
根据“无穷小 $\times$ 有界量 $=$ 无穷小”定理，极限为 0。
即 $\lim_{x \to 0} \frac{f(x)}{x} = 0$。
这意味着 $f(x)$ 是比 $x$ 高阶的无穷小。

\item \textbf{最终答案}
\textbf{答案}：D
\end{detailedanalysis}
\end{redenv}

% ========== 第4题 ==========
\problemrule
\section*{【题目4】确定无穷小阶数}
试确定当 $ x \rightarrow 0 $ 时，下列 \blank{1cm} 是关于 $x$ 的三阶无穷小。
\choicesTwo
{$ \sqrt[3]{x^{2}} - \sqrt{x} $}
{$ \sqrt{a + x^{3}} - \sqrt{a} 其中 a \neq 0 $}
{$ x^{3}+0.001x^{2} $}
{$ \sqrt[3]{\tan x^{3}} $}

\begin{redenv}
\analysisrule
\subsection*{逐步解析}

\begin{detailedanalysis}
\item \textbf{详细求解过程}
分别分析各选项：
\begin{enumerate}[label=\Alph*.]
\item $x^{2/3} - x^{1/2} \sim -x^{1/2}$（抓低阶）。阶数为 $1/2$。
\item 利用公式 $(1+\Delta)^\alpha - 1 \sim \alpha \Delta$。
$$ \sqrt{a+x^3} - \sqrt{a} = \sqrt{a}(\sqrt{1+\frac{x^3}{a}} - 1) \sim \sqrt{a} \cdot \frac{1}{2}\frac{x^3}{a} = \frac{1}{2\sqrt{a}} x^3 $$
这是 $x$ 的三阶无穷小（只要 $a \neq 0$）。
\item $x^3 + 0.001x^2 \sim 0.001x^2$。阶数为 2。
\item $\sqrt[3]{\tan x^3} \sim \sqrt[3]{x^3} = x$。阶数为 1。
\end{enumerate}

\item \textbf{最终答案}
\textbf{答案}：B
\end{detailedanalysis}
\end{redenv}

% ========== 第5题 ==========
\problemrule
\section*{【题目5】有界性判定}
函数 $ f(x)=\frac{|x|\sin(x-1)}{x(x-1)(x-2)} $ 在下列区间有界的是 \blank{1cm}。
\choicesFour{$ (0,1) $}{$ (1,2) $}{$ (0,2) $}{$ (2,3) $}

\begin{redenv}
\analysisrule
\subsection*{逐步解析}

\begin{detailedanalysis}
\item \textbf{问题本质分析}
函数在开区间有界，通常要求在区间的端点处的极限存在（或至少不为无穷大）。需检查 $x=0, 1, 2$ 处的极限。

\item \textbf{详细求解过程}
$$ f(x) = \frac{|x|\sin(x-1)}{x(x-1)(x-2)} $$
\begin{enumerate}[label=(\arabic*)]
\item \textbf{考察 $x \to 0$}：
$$ \lim_{x \to 0} \frac{|x|}{x} \cdot \frac{\sin(x-1)}{(x-1)(x-2)} = \text{有界} \cdot \frac{\sin(-1)}{(-1)(-2)} = \text{常数} $$
左极限为 -C，右极限为 +C。极限虽不存在但有限。所以在 $x=0$ 邻域有界。
\item \textbf{考察 $x \to 1$}：
$$ \lim_{x \to 1} \frac{\sin(x-1)}{x-1} \cdot \frac{|x|}{x(x-2)} = 1 \cdot \frac{1}{1 \cdot (-1)} = -1 $$
极限存在，有限。所以在 $x=1$ 邻域有界。
\item \textbf{考察 $x \to 2$}：
当 $x \to 2$ 时，分母 $(x-2) \to 0$，分子 $\to |2|\sin 1 \neq 0$。
极限为 $\infty$。所以在 $x=2$ 邻域无界。
\end{enumerate}

\textbf{分析选项}：
\begin{itemize}
    \item A. $(0, 1)$：左右端点极限均有限。有界。
    \item B. $(1, 2)$：右端点 $x=2$ 处无界。
    \item C. $(0, 2)$：包含 $x=2$ 的邻域，无界。
    \item D. $(2, 3)$：包含 $x=2$ 的邻域，无界。
\end{itemize}

\item \textbf{最终答案}
\textbf{答案}：A
\end{detailedanalysis}
\end{redenv}

% ========== 第6题 ==========
\problemrule
\section*{【题目6】寻找间断点}
函数 $ y=\frac{1}{e^{\frac{x}{x-2}}-1} $ 的间断点为 \blank{2cm}。

\begin{redenv}
\analysisrule
\subsection*{逐步解析}

\begin{detailedanalysis}
\item \textbf{详细求解过程}
间断点出现在分母为零或无定义的地方。
\begin{enumerate}[label=(\arabic*)]
\item \textbf{指数部分无定义}：$x-2=0 \implies x=2$。
\item \textbf{分母为零}：$e^{\frac{x}{x-2}} - 1 = 0 \implies e^{\frac{x}{x-2}} = 1 \implies \frac{x}{x-2} = 0 \implies x = 0$。
\end{enumerate}
所以间断点为 $x=0$ 和 $x=2$。

\item \textbf{最终答案}
\textbf{答案}：$ x=0, x=2 $
\end{detailedanalysis}
\end{redenv}

% ========== 第7题 ==========
\problemrule
\section*{【题目7】间断点分类}
$ f(x)=\begin{cases}\sin\frac{1}{x}, & x>0\\ x-1, & x\leq0\end{cases} $ 函数 $ f(x) $ 的间断点是第 \blank{2cm} 类间断点。

\begin{redenv}
\analysisrule
\subsection*{逐步解析}

\begin{detailedanalysis}
\item \textbf{详细求解过程}
考察分段点 $x=0$。
\begin{enumerate}[label=(\arabic*)]
\item \textbf{左极限}：$\lim_{x \to 0^-} f(x) = \lim_{x \to 0^-} (x-1) = -1$。
\item \textbf{右极限}：$\lim_{x \to 0^+} f(x) = \lim_{x \to 0^+} \sin\frac{1}{x}$。
这是一个振荡极限，值在 $[-1, 1]$ 之间变动，极限不存在。
\item \textbf{结论}：右极限不存在，属于第二类间断点。
\end{enumerate}

\item \textbf{最终答案}
\textbf{答案}：二（或振荡）
\end{detailedanalysis}
\end{redenv}

% ========== 第8题 ==========
\problemrule
\section*{【题目8】可去间断点}
$ x = 1 $ 为函数 $ y = \frac{x^{4}-1}{x^{3}-1} $ 的 \blank{2cm} 间断点。

\begin{redenv}
\analysisrule
\subsection*{逐步解析}

\begin{detailedanalysis}
\item \textbf{详细求解过程}
$x=1$ 处分母为0，无定义，是间断点。
计算极限：
$$ \lim_{x \to 1} \frac{x^4-1}{x^3-1} = \lim_{x \to 1} \frac{(x-1)(x+1)(x^2+1)}{(x-1)(x^2+x+1)} = \lim_{x \to 1} \frac{(x+1)(x^2+1)}{x^2+x+1} $$
$$ = \frac{2 \cdot 2}{3} = \frac{4}{3} $$
极限存在，故为可去间断点。

\item \textbf{最终答案}
\textbf{答案}：可去
\end{detailedanalysis}
\end{redenv}

% ========== 第9题 ==========
\problemrule
\section*{【题目9】可去间断点计数}
函数 $ f(x)=\frac{x^{3}-x}{\sin\pi x} $ 的可去间断点的个数为 \blank{2cm}。

\begin{redenv}
\analysisrule
\subsection*{逐步解析}

\begin{detailedanalysis}
\item \textbf{详细求解过程}
\begin{enumerate}[label=(\arabic*)]
\item \textbf{找间断点}：分母 $\sin\pi x = 0 \implies \pi x = k\pi \implies x = k \quad (k \in \mathbb{Z})$。
\item \textbf{分析可去性}：可去间断点要求极限存在，即分子也能抵消分母的零点。
分子 $x^3 - x = x(x-1)(x+1)$。
零点为 $x=0, 1, -1$。
\item \textbf{验证}：
\begin{itemize}
    \item $x=0$：$\lim_{x \to 0} \frac{x(x^2-1)}{\sin\pi x} = \lim_{x \to 0} \frac{x}{\sin\pi x} \cdot (-1) = \frac{1}{\pi} \cdot (-1)$。极限存在。
    \item $x=1$：$\lim_{x \to 1} \frac{x(x-1)(x+1)}{\sin\pi x}$。令 $t=x-1$，$\lim_{t \to 0} \frac{(t+1)t(t+2)}{\sin\pi(t+1)} = \lim \frac{1 \cdot t \cdot 2}{-\sin\pi t} = -\frac{2}{\pi}$。极限存在。
    \item $x=-1$：同理极限存在。
    \item $x=k (k \neq 0, 1, -1)$：分子不为0，极限为无穷大（无穷间断点）。
\end{itemize}
\item \textbf{结论}：可去间断点有3个（$0, 1, -1$）。
\end{enumerate}

\item \textbf{最终答案}
\textbf{答案}：3
\end{detailedanalysis}
\end{redenv}

% ========== 第10题 ==========
\problemrule
\section*{【题目10】连续性定义}
若要使 $ f(x)=\frac{1-\sqrt{1-x}}{1-\sqrt[3]{1-x}} $ 在点 $ x=0 $ 处连续，则要补充定义 $ f(0)= $ \blank{2cm}。

\begin{redenv}
\analysisrule
\subsection*{逐步解析}

\begin{detailedanalysis}
\item \textbf{此题关键}
连续 $\iff \lim_{x \to 0} f(x) = f(0)$。
即求极限。

\item \textbf{详细求解过程}
利用等价无穷小 $(1+x)^\alpha - 1 \sim \alpha x$。
令 $t = -x$，则 $x \to 0 \implies t \to 0$。
分子：$1 - (1+t)^{1/2} = -((1+t)^{1/2} - 1) \sim -\frac{1}{2}t = \frac{1}{2}x$。
分母：$1 - (1+t)^{1/3} = -((1+t)^{1/3} - 1) \sim -\frac{1}{3}t = \frac{1}{3}x$。
$$ \lim_{x \to 0} f(x) = \lim_{x \to 0} \frac{\frac{1}{2}x}{\frac{1}{3}x} = \frac{3}{2} $$

\item \textbf{最终答案}
\textbf{答案}：$ \frac{3}{2} $（或 1.5）
\end{detailedanalysis}
\end{redenv}

% ========== 第11题 ==========
\problemrule
\section*{【题目11】分段函数连续性}
设函数 $ f(x)=\begin{cases}x+a\cos2x, & x>0\\ e^{ax}-a, & x\leq0\end{cases} $ 为 $ (-\infty,+\infty) $ 上的连续函数，则 $ a= $ \blank{2cm}。

\begin{redenv}
\analysisrule
\subsection*{逐步解析}

\begin{detailedanalysis}
\item \textbf{详细求解过程}
在 $x=0$ 处必须连续，即左极限=右极限=函数值。
\begin{enumerate}[label=(\arabic*)]
\item \textbf{右极限}：$\lim_{x \to 0^+} (x + a\cos 2x) = 0 + a \cdot 1 = a$。
\item \textbf{左极限}：$\lim_{x \to 0^-} (e^{ax} - a) = e^0 - a = 1 - a$。
\item \textbf{函数值}：$f(0) = 1 - a$。
\item \textbf{联立方程}：$a = 1 - a \implies 2a = 1 \implies a = \frac{1}{2}$。
\end{enumerate}

\item \textbf{最终答案}
\textbf{答案}：$ \frac{1}{2} $
\end{detailedanalysis}
\end{redenv}

% ========== 第12题 ==========
\problemrule
\section*{【题目12】幂指函数极限}
求极限 $ \lim_{x\to+\infty}\left(x+\sqrt{1+x^{2}}\right)^{\frac{1}{x}}= $ \blank{2cm}。

\begin{redenv}
\analysisrule
\subsection*{逐步解析}

\begin{detailedanalysis}
\item \textbf{详细求解过程}
取对数求极限：
$$ L = \lim_{x \to +\infty} \frac{\ln(x + \sqrt{1+x^2})}{x} $$
这是 $\frac{\infty}{\infty}$ 型，使用洛必达法则：
$$ = \lim_{x \to +\infty} \frac{\frac{1}{x+\sqrt{1+x^2}}(1 + \frac{x}{\sqrt{1+x^2}})}{1} $$
$$ = \lim_{x \to +\infty} \frac{1}{x+\sqrt{1+x^2}} \cdot \frac{\sqrt{1+x^2}+x}{\sqrt{1+x^2}} $$
$$ = \lim_{x \to +\infty} \frac{1}{\sqrt{1+x^2}} = 0 $$
所以原极限为 $e^0 = 1$。

\item \textbf{最终答案}
\textbf{答案}：1
\end{detailedanalysis}
\end{redenv}

% ========== 第13题 ==========
\problemrule
\section*{【题目13】未知参数求解}
设函数 $ f(x)=\begin{cases}\frac{\sin2x+e^{2ax}-1}{x}, & x\neq0,\\ a, & x=0\end{cases} $ 在 $ (-\infty,+\infty) $ 上连续，则 $ a= $ \blank{2cm}。

\begin{redenv}
\analysisrule
\subsection*{逐步解析}

\begin{detailedanalysis}
\item \textbf{详细求解过程}
要求 $\lim_{x \to 0} f(x) = f(0) = a$。
计算极限：
$$ \lim_{x \to 0} \frac{\sin 2x + (e^{2ax} - 1)}{x} $$
利用 $\sin 2x \sim 2x$，$e^{2ax} - 1 \sim 2ax$。
$$ = \lim_{x \to 0} \frac{2x + 2ax}{x} = 2 + 2a $$
令 $2 + 2a = a$，解得 $a = -2$。

\item \textbf{最终答案}
\textbf{答案}：-2
\end{detailedanalysis}
\end{redenv}

% ========== 第14题 ==========
\problemrule
\section*{【题目14】零点定理应用}
证明方程 $ x^{5}+x-1=0 $ 只有一个正根。

\begin{redenv}
\analysisrule
\subsection*{逐步解析}

\begin{detailedanalysis}
\item \textbf{证明过程}
设 $f(x) = x^5 + x - 1$。
\begin{enumerate}[label=(\arabic*)]
\item \textbf{存在性}：
$f(0) = -1 < 0$；
$f(1) = 1 + 1 - 1 = 1 > 0$。
由于 $f(x)$ 连续，根据零点定理，在 $(0, 1)$ 内至少有一个根。
\item \textbf{唯一性}：
求导 $f'(x) = 5x^4 + 1$。
因为 $x^4 \ge 0$，所以 $f'(x) \ge 1 > 0$。
函数在 $\mathbb{R}$ 上严格单调递增。
所以方程只有一个实根（也就是只有一个正根）。
\end{enumerate}
\end{detailedanalysis}
\end{redenv}

% ========== 第15题 ==========
\problemrule
\section*{【题目15】连续性与参数}
设函数 $ f(x)=\begin{cases}x^{2}+1, & |x|\leq C\\ \frac{10}{|x|}, & |x|>C\end{cases} $ 在定义域上连续，试求常数 $C$。

\begin{redenv}
\analysisrule
\subsection*{逐步解析}

\begin{detailedanalysis}
\item \textbf{详细求解过程}
虽然有绝对值，但考虑到对称性，只需考虑 $x=C$ ($C>0$) 处的连续性。
\begin{enumerate}[label=(\arabic*)]
\item \textbf{函数值/左极限}：$f(C) = C^2 + 1$。
\item \textbf{右极限}：$\lim_{x \to C^+} \frac{10}{x} = \frac{10}{C}$。
\item \textbf{方程}：$C^2 + 1 = \frac{10}{C}$。
$$ C^3 + C - 10 = 0 $$
\item \textbf{求解}：试根法。
当 $C=2$ 时，$8 + 2 - 10 = 0$。符合。
检查单调性：$g(x) = x^3+x-10$，导数 $3x^2+1>0$，单调递增，故只有唯一实根。
\end{enumerate}

\item \textbf{最终答案}
\textbf{答案}：2
\end{detailedanalysis}
\end{redenv}

\end{document}
